\subsubsection{问题一具体分析}

问题一要求完成数据统计、未来 24 小时到达 GPU 需求预测与基础调度，属于"统计描述+预测+可行性优化"的综合类问题。本文采用"描述统计$ \to $时序特征工程$ \to $岭回归滚动回测预测$ \to $可行贪心构造"的思路，最后以统一能源结算口径核算目标值与 GPU 利用率。

首先对 50000 条任务按类型（AI 训练 16559 / 批量推理 16717 / 实时推理 16724）、区域、到达小时做统计，建立 6 区域$\times$3 类型的逐时到达 GPU 需求矩阵，并把有限 GPU 小时（占全部 0--2399 小时）识别为训练段；接着以"历史均值、$24/168$ 小时滚动均值、滞后项 $1/24/168$、昼夜谐波 $\sin/\cos$"构造共享特征，用岭回归在训练段滚动回测并以后向 24 小时为验证窗选型；然后以"实时任务到达即开工、弹性任务最早可行+最低时延区域"贪心构造调度；最后核算碳与成本，并输出最后 24 小时收尾甘特图。分析流程如图\ref{fig:q1_flow}所示。

\begin{figure}[H]
  \centering
  \begin{tikzpicture}[x=0.9cm, y=0.75cm]
    \node[box] (n1) at (0,0) {数据清洗};
    \node[box] (n2) at (4.5,0) {描述统计};
    \node[box] (n3) at (9,0) {特征工程};
    \node[box] (n4) at (13.5,0) {回测选型};
    \node[box] (n5) at (4.5,-2.4) {24h预测};
    \node[box] (n6) at (9,-2.4) {贪婪调度};
    \node[box] (n7) at (13.5,-2.4) {能源结算};
    \draw[arrow] (n1)--(n2); \draw[arrow] (n2)--(n3); \draw[arrow] (n3)--(n4);
    \draw[arrow] (n4)--++(0,-1.1)-|(n5); \draw[arrow] (n5)--(n6); \draw[arrow] (n6)--(n7);
  \end{tikzpicture}
  \caption{问题一分析流程}
  \label{fig:q1_flow}
\end{figure}

\subsubsection{问题一模型准备}

数据来自算力网络运行样例，任务表含 50000 条记录与任务 ID、来源区域、类型、到达小时、GPU 需求、时长、时延上限等字段；电力表含逐时电价、外送价、碳强度、可用/受限新能源、基线购售电与储能参数。检查发现的数据要点：六区域可用新能源逐时相同；GPU 需求以"等效 GPU 数"为口径，功率映射按类型给定；时延矩阵半开区间处理；"实时推理"任务时延上限为 $20$ ms 且要求到达即开工。预处理按类型合并功率映射、按区域合并 PUE/容量/储能参数，得到 6$\times$2407 的负荷—电力参数层与 50000 行的任务层。图\ref{fig:raw_q1}给出区域—类型逐时到达 GPU 需求；总到达 GPU·小时约 $5.02\times10^6$，而总容量小时约 $12.5\times10^6$，二者之比约 $40\%$，说明容量充裕、基础调度可行空间大。

\begin{figure}[H]
  \centering
  \includegraphics[width=0.92\textwidth]{figures/raw_q1_arrival_profile.png}
  \caption{六区域$\times$三类型小时到达 GPU 需求（0--2399）}
  \label{fig:raw_q1}
\end{figure}
